\section{Introduction}

Selecting an algorithm from instance features is a classical problem \cite{rice1976algorithm}.  Modern instance-space analysis turns features and performance into interpretable regions that reveal where algorithms succeed or fail \cite{smithmiles2009meta,smithmiles2023isa}.  The same perspective has reached quantum optimization and compilation.  It has been used to study parameter initialization \cite{katial2025qaoa}, quantum versus classical algorithm choice \cite{moussa2020selection}, compilation-option prediction \cite{quetschlich2023options}, and device-specific compiler selection \cite{quetschlich2025predictor}.  These studies establish that feature-conditioned performance mapping and algorithm selection are prior work, not contributions of this paper.

Exact compiler models permit a different question.  When an exact optimum or a complete proof is available, one can ask why a reference compiler fails, whether a restricted family always contains an optimum, and whether the feature representation itself determines the label.  These questions have different logical answers.  A support bound may be intrinsic or merely sufficient for one proof.  A proof can hold only for a region of objective weights.  A feature map can fit a complete finite domain yet fail on new sizes.  Combining these answers into one accuracy surface hides the distinction.

We therefore study a \emph{regime map} as a typed record rather than a single phase diagram.  The record separates six objects: the reference-exact region, constructive improvement mechanisms, intrinsic support, proof-derived support ceilings, objective-conditioned proof validity, and feature determination with prospective transfer.  A negative result is part of the map.  For example, mixed feature cells prove that no classifier using those features can be exact on the stated domain, while a failed prospective forecast shows that an in-domain rule did not transfer.

The evidence comes from four frozen quantum-compilation models (Table~\ref{tab:cross-model}).  Two Pauli block-encoding models admit all-size support results.  A six-term linear-combination model admits an exact incumbent boundary.  A weighted Clifford stabilizer-state model supplies complete finite graphs and a prospective size-transfer test.  The models are deliberately heterogeneous.  They are examples that separate logical coordinates, not a statistical sample of compilers.

The central result is that the coordinates do separate in practice.  In one model, an all-size support-two theorem coexists with a still-refinable vocabulary of named mechanisms.  In another, a safe syndrome-rank ceiling of five coexists with intrinsic support number one.  A proof-validity cone does not imply the opposite theorem outside the cone.  A near-injective state representation can determine every in-domain label and still fail out of sample.  These distinctions constrain how exact compiler results should be reported and compared.
